\section{Nuclear-recoil partition and charge yield}\label{app:E.23}

A silicon nucleus recoiling with \SI{100}{eV} cannot create an electron--hole pair in a binary collision with a valence electron, since that needs a recoil above the binary threshold $E_{R,\mathrm{th}} \approx \SI{352}{eV}$ of \cref{sec:19.2}. Yet the measured ionization yield at \SI{100}{eV} is $0.102^{+0.034}_{-0.030}$, and the data show no sign of an ionization threshold above \SI{100}{eV}~\cite[Table~II]{albakry2023first}. In germanium a recoil of \SI{254}{eV}, below its threshold of \SI{399}{eV}, gives \SI{64 \pm 8}{eV} of electron-equivalent energy, a quenching factor of $25\pm\SI{3}{\percent}$~\cite{kavner2024measurement}. The charge at these energies comes from processes that the Lindhard model of \cref{sec:19.3} does not contain, and the share of the recoil energy that stays in the lattice as stored energy is bounded by calibration data but not computed (\cref{sec:19.4}). Single recoils need the electronic excitation in a close collision of two atoms, the partition of the recoil energy into heat, stored energy and charge, and the charge yield below the binary threshold; dense tracks need the interaction of their carriers, the recombination of the track plasma and the amorphous pockets left by recoils of keV. Molecular dynamics on machine-learned potentials (\cref{app:E.11}) carries all of them except the excitation kernel and the track recombination. \Cref{tab:E23-summary} collects their grids, tests and costs.

\subsection{Close-collision excitation kernel}

Below $E_{R,\mathrm{th}}$, electrons are promoted at the level crossings of the quasi-molecule that two closely approaching atoms form, and \cref{sec:19.2} writes the result as the kernel of \cref{eq:19.2-kernel}, the distribution $\dd P/\dd\omega_e(b_{\mathrm{imp}}, E_{\mathrm{coll}})$ of the electronic excitation energy $\omega_e$ at impact parameter $b_{\mathrm{imp}}$ and relative energy $E_{\mathrm{coll}}$. The grid has 10 logarithmic values of $E_{\mathrm{coll}}$ from \SI{50}{eV} to \SI{5}{keV} and 8 values of $b_{\mathrm{imp}}$ from 0 to \SI{1.5}{\angstrom}, 80 points in all.

The Born--Oppenheimer order is not small here, so the calculation sits off the perturbative axes of $\mathfrak{T}$ (\cref{sec:8.5}): the nuclei are classical, the electrons are propagated in time. The Massey parameter $\xi_{\mathrm{M}}$ of \cref{sec:19.1} is 15 for a silicon atom of \SI{20}{eV}, so nonadiabatic excitation is suppressed as $\ee^{-15}$. At \SI{1}{keV}, $v_N = \SI{83}{km/s}$ and $\xi_{\mathrm{M}} \approx 2$, and it is not.

Real-time time-dependent density-functional theory with Ehrenfest nuclei propagates the collision, in the real-space codes Octopus~\cite{tancogne2020octopus} and SALMON (Scalable Ab initio Light--Matter simulator for Optics and Nanoscience) or the plane-wave code Qb@ll. Cells hold 216 atoms (\SI{16.3}{\angstrom} across) or more; the projectile is launched along random directions; and the excitation is counted by projecting the time-evolved orbitals on the conduction manifold of the ground state about \SI{20}{fs} after the collision, 20 to 50 trajectories per point. The time step is at most \SI{2}{as} near closest approach, the real-space grid at most \SI{0.2}{\angstrom} (or the equivalent plane-wave cutoff), and the pseudopotentials include semicore states and are hard enough for the close approach at $E_{\mathrm{coll}} \ge \SI{1}{keV}$. The cell has to be at least \SI{15}{\angstrom} across, so that the wake does not re-enter; the 216-atom cell on a \SI{0.2}{\angstrom} grid has $5.4\times10^5$ points, and its 432 occupied orbitals take about \SI{3.7}{GB}. A trajectory of about \SI{30}{fs} costs some 250 core-hours, and the 1600 to 4000 trajectories of the grid about $10^6$.

The adiabatic exchange--correlation functional misplaces level crossings that involve charge transfer, so the gap of the distorted configuration comes out too small. The theory uncertainty is therefore the spread between three treatments of a \SI{10}{\percent} subset of trajectories, which costs a few times $10^6$ core-hours more because the hybrid functional is some thirty times dearer per step: Ehrenfest with a hybrid functional (HSE), decoherence-corrected surface hopping~\cite{tully1990molecular}, and Landau--Zener probabilities from diabatic couplings of the Si$_2$ or Ge$_2$ quasi-molecule computed by multireference configuration interaction. The hybrid arm serves silicon only. In germanium HSE06 at the standard exchange fraction of one quarter inverts the $\mathrm{\Gamma}$ and L valleys and fails the band check of \cref{app:E.14} (valley order, both gaps, masses at L and $\mathrm{\Gamma}$), so a germanium hybrid arm would measure that defect together with the local-gap error; the germanium spread rests on the other two treatments until a functional passes the check. Ehrenfest dynamics gives the mean excitation, not its distribution, so the kernel is built from the projection statistics, never from the mean electronic energy. Above $E_{R,\mathrm{th}}$ the mean energy lost per unit path must approach the electronic stopping that the loss function gives (\cref{sec:19.3}). The kernel supplies the charge channel of slow recoils in \cref{sec:19.2} and the stochastic excitation events of the next two subsections.

\subsection{Partition tables}

\Cref{sec:19.1} splits the recoil energy into heat $E_{\mathrm{heat}}$, all of the phonon energy, stored energy $E_{\mathrm{st}}$ and charge, and \cref{sec:19.4} bounds the stored share with calibration data. The calculation tabulates the joint distribution $P(E_{\mathrm{heat}}, E_{\mathrm{st}}, N_{\mathrm{eh}}|E_R, \hat{\vect{r}}_R, T)$ for silicon and germanium on $E_R$ from \SI{1}{eV} to \SI{10}{keV} (40 logarithmic points), 100 recoil directions, and 100 initial conditions per direction, drawn from the zero-point and thermal distribution of the lattice. It also gives the threshold surface $E_{\mathrm{d}}(\hat{\vect{r}}_R)$ and the Frenkel-pair probability $p_{\mathrm{FP}}(E_R, \hat{\vect{r}}_R)$ of \cref{app:E.11}, and the spectrum of the phonons the recoil creates. The dynamics sits on $\mathfrak{T}$ where \cref{app:E.11} places it, at $n_B = 0$ with classical nuclei, and is controlled between the quantum--classical boundary at \SI{0.41}{eV} (\cref{sec:19.1}) and $E_{\mathrm{d}}$, 12.5 to \SI{36}{eV} by direction, where $\xi_{\mathrm{M}} \gg 1$; above, the close-collision kernel is attached as stochastic events.

The potential is a machine-learned interatomic potential (MLIP) of the fast linear atomic cluster expansion (ACE) family, spliced to the screened Coulomb repulsion as in \cref{app:E.11}. It is trained on about $10^4$ configurations computed with PBE and SCAN: dimers down to \SI{0.5}{\angstrom}, point defects, close pairs, amorphous and liquid structures, and snapshots of recoils. It is accepted at force errors below \SI{50}{meV/\angstrom} and energy errors below \SI{3}{meV} per atom. $E_{\mathrm{d}}$ and $p_{\mathrm{FP}}$ follow the protocol of \cref{app:E.11}, its directions, energy steps, cell, border and run length, with this potential. What differs is the quench at $t_{\mathrm{pr}}$ below and a check of twenty directions by Born--Oppenheimer molecular dynamics in density-functional theory, 512 atoms, about $10^5$ core-hours. The partition tables use cells from $10^3$ atoms below \SI{100}{eV} to $10^5$ atoms at \SI{10}{keV}, run for 1 to \SI{10}{ps}, and cost $10^4$ to $10^5$ hours on graphics processors per material.

The stored energy is defined by the quench of \cref{app:E.11}, \cref{eq:E.11-Est}, at a time $t_{\mathrm{pr}}$ of 10 to \SI{100}{ps} after the recoil, by steepest descent or the fast inertial relaxation engine (FIRE)~\cite{bitzek2006structural} at fixed cell, without annealing. Its reference is the relaxed perfect cell, computed once for each potential, cell and boundary condition, and never the initial configuration of the run. The difference is large. The initial atoms are drawn from the zero-point distribution and carry half of the zero-point energy as potential energy. \Cref{sec:5.3} gives a zero-point energy of \SI{123}{meV} per cell of silicon, so the initial configuration lies \SI{31}{meV} per atom, or \SI{126}{eV} in a 4096-atom cell, above the minimum, 17 to 19 times the 6.8 to \SI{7.5}{eV} that a close Frenkel pair stores (\cref{sec:19.1}). The change of zero-point energy between the end basin and the perfect crystal comes from the Hessians at both minima; it is booked in the heat, not in $E_{\mathrm{st}}$. One table of configurational energies, with one potential, cell and quench, serves the stored energy here, the end states of the release paths of \cref{app:E.22} and the pockets below; otherwise the end states and the stored energies disagree by the potential's error on each defect. The spread of its entries between the potential and density-functional theory is the potential's uncertainty on $E_{\mathrm{st}}$.

The created-phonon spectrum is recorded at $t_{\mathrm{pr}}$, before the quench. The excess lattice state is projected on the harmonic modes $s$ of the quenched end state, $E_s = \tfrac12(p_s^2 + \omega_s^2u_s^2)$, and the same projection of the initial state is subtracted to remove the zero-point spectrum. The result is binned at \SI{0.5}{meV}, with branch labels from the overlap of eigenvectors with those of the perfect crystal. The energy that the projection misses, the anharmonic residual, must stay below \SI{10}{\percent}. The energy absorbed by the boundary thermostat is recorded with its spectrum, from the flux into the thermostat layer, or the spectrum is taken before the front reaches it. The acceptance test for the spectrum is the hand-over of \cref{sec:19.1}: for recoils small enough to stay harmonic, the created minus absorbed spectrum must reproduce the phonon spectrum of the nuclear response of \cref{sec:11.3}. The tables give the partition of \cref{sec:19.1} and the computed stored energy that \cref{sec:19.4} sets against its bounds.

\subsection{Charge yield below the binary threshold}

The measured silicon yields at 100 to \SI{460}{eV} straddle $E_{R,\mathrm{th}}$, and no calculation yet predicts them (\cref{sec:19.2}). Predicting them needs the full $P(N_{\mathrm{eh}}|E_R, \hat{\vect{r}}_R, T)$ for $N_{\mathrm{eh}} \le 30$ on $E_R$ from \SI{20}{eV} to \SI{1}{keV}, with at least 25 logarithmic points, dense near 352 and \SI{399}{eV}. From it follow the yield $Y_q(E_R) = \langle N_{\mathrm{eh}}\rangle\varepsilon_{\mathrm{eh}}/E_R$, the probability $P(N_{\mathrm{eh}} = 0|E_R)$ of a recoil without charge, the Fano factor $F_{\mathrm{NR}}(E_R)$ of the nuclear-recoil charge, and the correlations of $N_{\mathrm{eh}}$ with the heat and with $E_{\mathrm{st}}$ at fixed $E_R$, which \cref{sec:19.5} needs.

The trajectories are those of the partition tables, with excitation events drawn from the close-collision kernel at every close approach and the electronic friction of the potential switched off where the kernel acts. Three channels are flagged separately: binary excitation above $E_{R,\mathrm{th}}$, Landau--Zener promotion, and linear-response excitation through the high-momentum tail of the valence density, which makes the binary threshold soft (\cref{sec:19.2}). Each excitation $\omega_e$ is converted into pairs with the cascade distribution $P_{\mathrm{casc}}$ at that seed (\cref{app:E.13,app:E.24}). Slow-ion excitations lie near the gap, where the cascade is far from its asymptote, so dividing $\omega_e$ by the high-energy $\varepsilon_{\mathrm{eh}}$ is wrong. The Fano factor follows from the law of total variance, $\operatorname{Var}N_{\mathrm{eh}} = \mathbb{E}[\operatorname{Var}(N_{\mathrm{eh}}|\text{seeds})] + \operatorname{Var}(\mathbb{E}[N_{\mathrm{eh}}|\text{seeds}])$, which separates the fluctuations of the cascade from those of the excitation events.

The number of trajectories per point is set by the variance. For a sample of $n$ values with excess kurtosis $\kappa$, the relative error of the sample variance is $[2/(n - 1) + \kappa/n]^{1/2}$; for \SI{6}{\percent} at $\kappa = 2$ this needs $n \ge 1112$. The cost lies within the molecular-dynamics budget of the partition tables, about $10^3$ to $5\times10^3$ hours on graphics processors. The result is validated, not fitted, against the silicon data: $Y_q = 0.102^{+0.034}_{-0.030}$ and $F_{\mathrm{NR}} = 0.9^{+0.7}_{-0.4}$ at \SI{0.10}{keV}, $0.108^{+0.009}_{-0.010}$ and $0.5^{+0.2}_{-0.1}$ at \SI{0.22}{keV}, $0.136^{+0.009}_{-0.008}$ and $1.8^{+0.6}_{-0.5}$ at \SI{0.46}{keV}~\cite[Table~II]{albakry2023first}, and against the germanium point at \SI{254}{eV}~\cite{kavner2024measurement}. The silicon yields assume \SI{3.8}{eV} per pair and a total phonon energy equal to the recoil energy plus the work of the pairs in the bias field, so that no energy is stored~\cite[Eq.~(1)]{albakry2023first}. The computed yield is converted to the same definition before the comparison. The stored energy, at most a few percent of $E_R$ (\cref{sec:19.4}), lies inside their errors.

\subsection{Dense tracks}

The cascade of \cref{sec:18.1} treats the pairs of one event as independent, each subtree evolving alone. At keV, a nuclear-recoil track packs its carriers so densely that they interact: the independence is the zeroth order of an expansion in the pair density times the exciton volume, $\nd_{\mathrm{pair}}a_x^3$, and the first correction is a carrier--carrier collision term. In germanium the lattice itself can melt locally in a thermal spike, beyond the description of the ion as a sequence of binary collisions. The calculation gives the ratio of $P(N_{\mathrm{eh}}|E_R)$ with and without interactions and the change of $F_{\mathrm{NR}}$, per recoil energy and species, which \cref{sec:19.1} uses to bound the independence of the subtrees.

The carriers are followed by ensemble Monte Carlo with molecular-dynamics Coulomb forces among all of them, computed by the particle--particle particle--mesh method, seeded with the end-of-cascade distribution and the geometry of the track, with screened impact ionization and Auger recombination as extra channels. The lattice is followed by machine-learned molecular dynamics of the recoil coupled to a two-temperature model of the electrons~\cite{duffy2007including}, with the electron--phonon coupling from the vertex of \cref{app:E.5}. The time step is \SI{0.1}{fs}, the Coulomb mesh at most \SI{1}{nm}, the electron-temperature grid at most \SI{1}{nm}, and at least $10^3$ tracks are run per energy and species. The dynamics costs about $10^4$ hours on graphics processors and the carriers $10^3$ core-hours. In the dilute limit, with the carrier density scaled down, the calculation must recover the independent cascade. Three errors must be avoided: the electronic stopping friction of the potential and the close-collision kernel must not both act; carriers that leave the molecular-dynamics box must be continued in the cascade, not dropped; and the thermostat must not absorb the heat of the spike before it has spread.

\subsection{Track recombination}

A dense track is a plasma of electrons and holes that shields an applied field from its interior and loses carriers by Auger recombination, at a rate $C_{\mathrm{Aug}}\nd^3$ per volume. A first estimate shows which parameters matter. Take a cylindrical Gaussian track of $\lambda_{\mathrm{tr}}$ pairs per length and initial radius $r_{\mathrm{tr}}$, spreading by ambipolar diffusion with constant $D_{\mathrm{amb}}$ at zero field. Its density is $\nd = \lambda_{\mathrm{tr}}\ee^{-\rho^2/s}/\pi s$ with $s = r_{\mathrm{tr}}^2 + 4D_{\mathrm{amb}}t$, and $\int\dd^2\rho\,\nd^3 = \lambda_{\mathrm{tr}}^3/3\pi^2s^2$. To first order in $C_{\mathrm{Aug}}$ the fraction lost is
\begin{equation}
  f_{\mathrm{A}} = \frac{C_{\mathrm{Aug}}\lambda_{\mathrm{tr}}^2}{3\pi^2}\int_0^\infty\frac{\dd t}{(r_{\mathrm{tr}}^2 + 4D_{\mathrm{amb}}t)^2}
  = \frac{C_{\mathrm{Aug}}\lambda_{\mathrm{tr}}^2}{12\pi^2D_{\mathrm{amb}}r_{\mathrm{tr}}^2} .
  \label{eq:E.23-fA}
\end{equation}
The loss is set at early times, as $1/r_{\mathrm{tr}}^2$, so the initial radius from the dense-track and pocket calculations matters more than the late diffusion. In a biased detector the plasma screens the field until diffusion erodes it; the time this takes, the plasma time, delays the induced current.

The full calculation solves drift--diffusion and Poisson equations with rate terms for Auger recombination, exciton formation and Shockley--Read--Hall capture on the geometry of the track, two-dimensional with axial symmetry, with a mesh of at most \SI{1}{nm} near the core and fields from 0 to \SI{250}{V/cm}. It gives $f_{\mathrm{A}}(E_R, \mathcal{E}_{\mathrm{b}})$, with $\mathcal{E}_{\mathrm{b}}$ the bias field, the plasma time and the delay of the induced current, for under $10^3$ core-hours, and it must reproduce \cref{eq:E.23-fA} for $f_{\mathrm{A}} \ll 0.1$; above $f_{\mathrm{A}} \approx 0.1$ the loss is nonlinear and only the full solution holds. The weak input is $C_{\mathrm{Aug}}$, which has not been measured at millikelvin in silicon or germanium; it is carried as a stated parameter, scanned from a tenth to ten times its room-temperature value. Exciton formation in the track does not by itself remove gap energy, since a bound pair still holds $\Eg - E_x^b$; only loss channels, Auger recombination and exciton--exciton annihilation, do. The computed loss and its field dependence feed the dense-track charge of \cref{sec:19.1} and the kinetics of \cref{app:E.25}.

\subsection{Amorphous pockets}

Recoils of keV in silicon and germanium leave small amorphous pockets as well as point defects. A pocket that recrystallizes later releases its stored energy at once, and \cref{sec:19.6} counts such releases among the candidate sources of the low-energy excess. A release is visible only if its energy, after collection, exceeds the threshold; with a recrystallization energy of order \SI{0.1}{eV} per atom, an estimate, a threshold of tens of eV needs pockets of hundreds of atoms. The calculation gives the size distribution $P(N_a|E_R, \hat{\vect{r}}_R)$ of the number $N_a$ of atoms in pockets at $t_{\mathrm{pr}}$, and the spectrum of recrystallization barriers.

The calculation sits at the same place on $\mathfrak{T}$ as the partition tables. The cascades are run with the spliced machine-learned potential at $E_R$ from 0.2 to \SI{20}{keV}, at least $10^3$ per energy and material, in cells of $10^5$ to $10^6$ atoms for 10 to \SI{50}{ps} with a boundary thermostat, and quenched as above. A pocket is a connected set of atoms whose local order parameter, a bond-angle distribution or a Steinhardt parameter, falls below a crystalline cut. The barrier spectrum comes from nudged elastic bands and accelerated molecular dynamics on about 100 pockets. The cost is $10^4$ to $10^5$ hours on graphics processors. The definition of a pocket changes $N_a$ by a factor of about two, so the distribution is reported for two cuts. Thermal spikes in germanium need the two-temperature coupling of the dense-track calculation. The pocket energies come from the configurational table above, so that pockets and point defects add up to the stored energy of the partition tables. The size distribution and barriers give the pocket term of the delayed-release spectrum of \cref{sec:19.6}.

\begin{table}[t]
  \caption{Nuclear-recoil calculations for silicon and germanium. Costs are estimates. GPU-h are hours on graphics processors; TDDFT, time-dependent density-functional theory; MLIP-MD, molecular dynamics with a machine-learned interatomic potential; MC, Monte Carlo.}
  \label{tab:E23-summary}
  \centering
  \small
  \begin{tabular}{p{0.19\linewidth}p{0.29\linewidth}p{0.27\linewidth}p{0.13\linewidth}}
    \toprule
    quantity & method and grid & acceptance & cost \\
    \midrule
    $\dd P/\dd\omega_e(b_{\mathrm{imp}}, E_{\mathrm{coll}})$ & real-time TDDFT on 216 atoms, step \SI{2}{as}, grid \SI{0.2}{\angstrom}; $10\times8$ points from \SI{50}{eV} to \SI{5}{keV} & stopping above $E_{R,\mathrm{th}}$; HSE (Si only), surface hopping, Landau--Zener spread & $10^6$ core-h; hybrid subset a few $10^6$ \\
    $P(E_{\mathrm{heat}}, E_{\mathrm{st}}, N_{\mathrm{eh}})$, $E_{\mathrm{d}}$, $p_{\mathrm{FP}}$, phonon spectra & MLIP-MD, 40 energies $\times$ 100 directions $\times$ 100 seeds & DFT on 20 directions; harmonic hand-over & $10^4$--$10^5$ GPU-h \\
    $Y_q$, $F_{\mathrm{NR}}$, $P(N_{\mathrm{eh}} = 0)$ & MD with excitation events, \SI{20}{eV}--\SI{1}{keV}, $\ge 1112$ runs per point & silicon data at 0.10, 0.22, \SI{0.46}{keV} & $10^3$--$5\times10^3$ GPU-h \\
    Dense-track ratio, $\Delta F_{\mathrm{NR}}$ & carrier MC with Coulomb forces; MD with two-temperature model & dilute limit & $10^4$ GPU-h \\
    $f_{\mathrm{A}}$, plasma time & drift--diffusion--Poisson, mesh $\le\SI{1}{nm}$ & \cref{eq:E.23-fA} for $f_{\mathrm{A}} \ll 0.1$ & $10^3$ core-h \\
    $P(N_a|E_R)$, barrier spectrum & MLIP-MD, 0.2--\SI{20}{keV}, $10^5$--$10^6$ atoms & two cluster cuts; energy closure & $10^4$--$10^5$ GPU-h \\
    \bottomrule
  \end{tabular}
\end{table}
